% =====================================================================
\section{What the missing tool must satisfy}

\figlean
\label{sec:missing}
% =====================================================================

The campaign's sharpest deliverable is a specification of the object it could not build. The
species must be \textbf{one} tool satisfying \textbf{five} conditions, stated verbatim from
\code{docs/PROOF-SKELETON.md} \S3 and \code{docs/new-tools-registry.md}:

\begin{enumerate}
\item \textbf{Nonperturbative.} It must be valid outside every convergent expansion, at
  $\beta\approx2.5$, where the relevant $g\approx1.5$ and no small parameter exists. (A
  perturbative resummation is excluded by the Savvidy negative mode and the IR renormalon/Borel
  ambiguity of the $a^{-4}$ coefficient.)
\item \textbf{Sign-controlling lower bound.} It must produce a \emph{lower} bound
  $m_{\mathrm{eff}}^2(k\to0)\ge\lambda>0$ (equivalently $G''\ge\lambda$, $f_1^2\ge r_0$), not
  another upper bound. This is the rarest direction: almost every physical tool gives upper
  bounds or sign constraints.
\item \textbf{Non-circular.} It must not assume the gap, the threshold, or the isolation it is
  meant to prove. (Temple/Weinstein, the BS threshold, the second-level $E_2$, and the
  cutoff-independent Poincar\'e floor all fail here by presupposing the answer.)
\item \textbf{Non-SUSY and fails the free field.} It must not rely on supersymmetry (the physical
  theory is pure YM) and it must \textbf{fail} on the free massless field, where $f_1^2=0$,
  $G''=0$. This is the acid test: a tool that ``proves'' a gap for the free field is vacuous.
\item \textbf{Uniform lower bound.} It must be uniform in $(L,a)$ --- volume- and
  cutoff-uniform --- and survive the continuum limit. A fixed-$a$ or fixed-$L$ statement does not
  give Clay.
\end{enumerate}

\subsection{The equivalent forms of the object}

These five conditions define the same object in several vocabularies
(\code{docs/YANG-MILLS-MASTER.md} \S6):

\begin{enumerate}
\item \textbf{H1:} coarse Poincar\'e $\operatorname{Var}_\nu(\varphi)\le(1/\lambda)\int|\nabla
  \varphi|^2$, uniform $\lambda>0$.
\item \textbf{(CP):} a single $T=0$ confining phase (no bulk transition) $+$ a $0^{++}$ edge that
  is Kato-preserved $+$ uniform isolation.
\item \textbf{Edge atom $f_1^2>0$:} $\nu(\{m_1^2\})>0$ (usc, non-diagonal, needs $E_2$).
\item \textbf{$G''(0)>0$ at $k\to0$:} the coarse curvature, $=S-D$; the reported $D/S\approx2\%$
  is on the gauge-dependent gluon mode, not the $0^{++}$ (Appendix~\ref{app:postaudit}).
\item \textbf{Determinant zero in the continuum:} $\det(1-zT_O)$ at $a\to0$.
\item \textbf{Savvidy sign} (non-abelian): the indefinite Hessian's sign resolved.
\end{enumerate}

These are not six problems; they are one --- the volume-uniform lower bound on $m_{0^{++}}$ along
the path, which is the mass gap itself. The final convergence reduces them to \textbf{one cell}:
the ``typical field $\times$ high momentum'' cell of the $(\text{field}\times\text{momentum})$
plane, requiring the ceiling $\operatorname{Var}(\xi_{k\to0})\le1/\lambda$ (curvature form
$\operatorname{Var}_\mu(\langle v,\nabla S\rangle|0)\le E_\mu[H_{vv}|0]+W_{\mathrm{full}}-\delta$;
Appendix~\ref{app:sharp}). Every other cell is closed: low field by $\Phi''\succeq N/12+$ BL, the
high-field cone by WSTC (measure starvation), the AF-UV cone by $g\to0$. \textbf{That cell is a
variance-upper-bound species, the same species as G1's L3, at a different scale.}

\subsection{The sharpest form, with no slack}

The single open inequality, stated with no slack, is: there is $\Delta_0>0$ such that, uniformly
in $(L,a)$ and for every coarse transverse $v$ at $k\to0$,
\begin{equation}
\operatorname{Var}_\mu\!\big(\langle v,\nabla V\rangle\,\big|\,\xi=0\big)
\;\le\;
E_\mu\!\big[H_{vv}\,\big|\,\xi=0\big]\;-\;a^{2}\Delta_0^{2},
\qquad V=S_W+\Phi .
\label{eq:sharpest}
\end{equation}
This is the (M) form with $W_{\mathrm{full}}=W_{\mathrm{geo}}+X$ and
$\delta:=a^{2}\Delta_0^{2}+W_{\mathrm{full}}$; the two displays are the same inequality.
Equivalently $D\le S-\lambda$ \emph{together with} $S\ge\lambda$ (the bare defect form
$D\le(1-\epsilon)S$ is insufficient), equivalently the dimensionless ratio
$G''_a/a^{2}\ge\Delta_0^{2}$.

Two dimension traps are recorded. The unconditioned budget $E[f^{2}]\approx2$ is
ultraviolet/equipartition-pinned and \emph{not} proportional to $a^{2}$, so it cannot stand in for
$\lambda$; and the sign lives entirely in the \emph{conditioned} fibre Hessian
$K=\Pi_0\operatorname{Hess}V\,\Pi_0$, whose negative edge ($\det K=0$, the transverse massless
mode) is the wall. The unconditioned expected Hessian is provably PSD,
$E[\operatorname{Hess}V]=E[\nabla V\otimes\nabla V]\succeq0$; the gap is therefore the
\emph{conditioning correction} --- $W_{\mathrm{geo}}+X$ together with the change in slice variance
--- i.e.\ the step from the expected curvature to the conditioned one.

\subsection{The machine that does not exist}

To close Clay one must build \textbf{a nonperturbative machine that transports the strong-coupling
gap and residue continuously across the Savvidy crossing and takes $a\to0$}. The campaign states
the verdict plainly (\code{docs/yang-mills-frontier.md} \S7, \code{YANG-MILLS-MASTER.md} \S9):

\begin{qbox}
The last step is to give a volume-uniform lower bound on $m_{0^{++}}$ across $[1/48,2.5]$. Every
tried tool stops at the strong-coupling end (SZZ, large-$N$) or on the wrong half (UV). What is
needed is a machine that carries the strong-coupling gap across the Savvidy/sign crossover
continuously to the physical point --- \textbf{this machine does not exist.} It is equivalent to
nonperturbative 4D dynamics, the classification of nonperturbative pure-gauge 4D CFTs
(\code{TOOL-not-iii.md}).
\end{qbox}

\subsection{Why the existing rigorous tools stop}

It is worth saying precisely why the modern rigorous toolbox does not reach the crossover, since
the statement is often met with disbelief.

\begin{itemize}
\item \textbf{Balaban-type renormalization group} controls the ultraviolet and reaches
  $\lambda\sim1$. It is a small-field, multiscale expansion; the crossover is at
  $\lambda_{\mathrm{tH}}\in[1,7.3]$, where the expansion parameter is $O(1)$ and no small field
  is available.
\item \textbf{Strong-coupling expansions} (SZZ, random surfaces, center vortices) require
  $\lambda$ large. They reach $\beta<1/48$ or $\lambda_{\mathrm{tH}}<1/24$; the physical coupling
  is far outside.
\item \textbf{Constructive QFT with fermions} (Magnen--Rivasseau--S\'en\'eor and successors)
  constructs the theory with an infrared cutoff; the mass gap is not uniform in the cutoff.
\item \textbf{Rigorous large-$N$} (Chatterjee) is a strong-coupling statement about the lattice
  measure; it does not reach the physical coupling.
\item \textbf{Stochastic quantization / regularity structures} (Hairer and collaborators) are
  small-parameter theories; 4D Yang--Mills is not subcritical in the required sense, and the
  large-field (Savvidy-indefinite) estimate is the missing ingredient.
\item \textbf{Validated numerics} control finite-volume, finite-cutoff operators; they give upper
  bounds on the continuum gap, not lower bounds.
\end{itemize}
The strongest contemporary rigorous frontier (CPS/CNS 2025) reaches only
$\lambda_{\mathrm{tH}}<1/24$ (the older ``$75\times$'' figure was a unit confusion; see
\code{PROOF-EMERGENT-4D.md}), nowhere near $\beta\approx2.5$.

\subsection{What building it would close}

A single one-sided variance upper bound, nonperturbative, coupling-free, $(L,a)$-uniform, failing
at $f_1^2=0$, would close \textbf{both} G1's L3 (UV) and G2's H1 (IR) simultaneously. That is the
precise sense in which ``one species, two scales'' is the whole gap between the campaign and
Clay. In the $\beta$ coordinate, the same object is a nonperturbative sign theorem for
$\beta_{N_f=0}$ or a transport across the Savvidy crossing; in the measure coordinate it is the
negative-association inequality or an equivalent discontinuous functional. The three descriptions
are the same machine.

\subsection{The three named open inequalities}

For the reader who wants the shortest possible statement of what is missing, here are the three
inequalities, each named, each with no slack, each equivalent to the gap in its coordinate.

\begin{enumerate}
\item[\textbf{(H1)}] $\displaystyle
  \operatorname{Var}_\mu\bigl(\langle v,\nabla S\rangle\mid0\bigr)
  \le E_\mu[H_{vv}\mid0]+W_{\mathrm{full}}-\delta$, uniformly in $(L,a)$, every coarse $v$ at
  $k\to0$.
\item[\textbf{(NA)}] $\displaystyle \operatorname{Cov}_\mu\bigl(G,\Phi'\mid0\bigr)\le0$ for the
  frustrated $\mathrm{SU}(N)$ Haar $\times$ Wilson measure --- the only variance-free face.
\item[\textbf{(TCI)}] $\displaystyle
  \langle\varphi,Q\varphi\rangle+\langle\varphi,\mathbb{K}\varphi\rangle-cF\ge\mu>0$ for all
  unit $\varphi\in\mathcal{H}_-$, pointwise and nonperturbatively --- the tachyonic-compensation
  inequality, equivalently convexity of the gauge-invariant effective potential at its minimum.
\end{enumerate}
Each is a lower bound on a curvature or an upper bound on a variance; each fails on the free
field; each is nonperturbative; each is $(L,a)$-uniform in its statement. Proving any one of them
would close the gap. None is proved.

\subsection{Honest bottom line}

The specification is complete. The object is named, its five conditions are stated, its
equivalent forms are displayed, and its failure modes are diagnosed. What is missing is the
object itself. This is the campaign's real result: not a proof, but a precise description of the
shape of the hole. The next chapter gives the verdict.
